\documentclass[aps,reprint,amsmath,amssymb,nofootinbib,floatfix]{revtex4-2}

\usepackage{graphicx}
\usepackage{microtype}
\usepackage[hidelinks]{hyperref}
\usepackage{float}

\graphicspath{{../../outputs/yang_fixation_circular_working_memory/perturbation_experiments/noise_structure/gaussian_control/figures/}}

\begin{document}

\section{Independent Gaussian noise as a control perturbation}

\subsection{Rationale}

Psilocybin has been associated with increased neural signal diversity and has
motivated accounts of psychedelic states in terms of increased entropy or
variability~\cite{carhartharris2014,schartner2017}. These findings motivate
testing whether a working-memory network is sensitive to stochastic
perturbation, but they do not imply that psilocybin acts as additive Gaussian
noise. We therefore used independent Gaussian noise as a deliberately
non-specific control rather than as a mechanistic model of the drug. This
choice also has a direct computational rationale: stochastic fluctuations can
drive diffusion of maintained representations, and the extent of this drift
depends on the stability and geometry of the recurrent memory
dynamics~\cite{kilpatrick2013,ghazizadeh2021}. The control consequently asks a
restricted question: how much degradation follows from adding unstructured
variability of a known magnitude to an otherwise unchanged trained network?

\subsection{Implementation}

Noise was added to the hidden-state update only during the memory delay. For
each active time step $t$, trial $b$, and hidden unit $i$, an independent sample
$\eta_{tbi}\sim\mathcal{N}(0,1)$ was drawn. The realised perturbation was

\begin{equation}
\varepsilon_{tbi}=
\sigma\frac{\eta_{tbi}}
{\left(|\mathcal{A}|^{-1}\sum_{(t,b,i)\in\mathcal{A}}
\eta_{tbi}^{2}\right)^{1/2}},
\label{eq:gaussian-control}
\end{equation}

where $\mathcal{A}$ denotes all delay-active entries and $\sigma$ is the target
root-mean-square (RMS) perturbation strength. This normalisation matched the
realised RMS exactly within each evaluated batch. Noise was independent across
time, trials, and units; no temporal autocorrelation or population covariance
was imposed. Five independently trained frozen models were evaluated at delay
lengths of 20, 80, and 160 steps and at RMS strengths $0$, $0.01$, $0.025$,
$0.05$, and $0.10$. Each model--delay condition comprised five stochastic
replicates of 20 paired 64-trial batches. A linear decoder fitted to 512 clean
trials per model and delay quantified maintenance error, and the trained model
seed was treated as the replication unit ($n=5$). RMS is an intervention
parameter and should not be interpreted as a pharmacological dose.

\subsection{Results}

Gaussian noise produced a graded, delay-dependent impairment
(Fig.~\ref{fig:gaussian-behaviour}). At RMS $0.05$, mean response error increased
from $4.34\pm1.54^{\circ}$ to $4.91\pm1.45^{\circ}$ at delay 20, from
$6.49\pm1.93^{\circ}$ to $7.47\pm1.68^{\circ}$ at delay 80, and from
$9.03\pm2.50^{\circ}$ to $10.44\pm2.24^{\circ}$ at delay 160 (mean $\pm$ SD
across five model seeds). The impairment became clearer at RMS $0.10$, where
the corresponding response errors were $6.42\pm1.35^{\circ}$,
$10.48\pm1.92^{\circ}$, and $14.68\pm2.57^{\circ}$.

\begin{figure}[H]
\centering
\includegraphics[width=\linewidth]{figure_g1_dose_delay_behaviour.pdf}
\caption{\textbf{Behavioural response to independent Gaussian perturbation.}
Absolute response error is shown across perturbation strengths and delay
lengths. Solid lines are five-model means; dashed lines and shading denote 95\%
bootstrap confidence intervals across models.}
\label{fig:gaussian-behaviour}
\end{figure}

The maintenance measures showed a larger relative effect than the final
behavioural readout (Fig.~\ref{fig:gaussian-maintenance}). At delay 80 and RMS
$0.05$, clean-decoder error increased from $1.29\pm0.82^{\circ}$ to
$3.71\pm1.53^{\circ}$, while decoded angular drift increased from
$2.60\pm1.84^{\circ}$ to $5.20\pm2.11^{\circ}$. Fixation/wait--go accuracy was
effectively preserved ($0.98\pm0.00$ in both conditions), indicating that the
perturbation primarily affected the maintained angle rather than task-phase
gating. The decoder was fitted to clean hidden states and then evaluated on
the corresponding perturbed states, so clean-decoder error measures
representational distortion rather than a refitted noisy readout.

\begin{figure}[H]
\centering
\includegraphics[width=\linewidth]{figure_g2_maintenance_metrics.pdf}
\caption{\textbf{Maintenance effects at delay 80.} Clean-decoder error and
decoded drift increased with Gaussian-noise RMS, whereas fixation/wait--go
accuracy remained stable. Lines and intervals are defined as in
Fig.~\ref{fig:gaussian-behaviour}.}
\label{fig:gaussian-maintenance}
\end{figure}

The geometry analysis was consistent with mild blurring rather than collapse
of the circular representation (Fig.~\ref{fig:gaussian-geometry}). At delay 80
and RMS $0.05$, within-angle dispersion increased from $0.550$ to $0.732$,
whereas normalised between-angle centroid spacing was $0.990$ relative to the
clean value of $1$. Thus, trial-to-trial states for the same remembered angle
became more dispersed while the global separation of different angles was
largely retained.

\begin{figure}[H]
\centering
\includegraphics[width=\linewidth]{figure_g3_geometry.pdf}
\caption{\textbf{Overlay of clean and Gaussian memory geometry at delay 80 and
RMS 0.05.} Angle centroids and uncertainty ellipses are shown in the
clean-fitted PC1--PC2 plane. The Gaussian ring is overlaid on the clean
baseline so that its mild within-angle broadening and small centroid changes
can be read directly.}
\label{fig:gaussian-geometry}
\end{figure}

Response settling was close to its measurement floor at the primary RMS
(Fig.~\ref{fig:gaussian-settling}). At delay 80, the mean first response step
within $20^{\circ}$ of the target increased only from $5.02$ in the clean
condition to $5.19$ at RMS $0.05$; $99.2\%$ of trials settled. At RMS $0.10$,
settling increased to $6.55$ steps and the settled fraction fell to $93.6\%$.
Final response error in this focused analysis increased from $5.94^{\circ}$
clean to $6.75^{\circ}$ at RMS $0.05$ and $10.40^{\circ}$ at RMS $0.10$.

\begin{figure}[H]
\centering
\includegraphics[width=\linewidth]{figure_g4_settling.pdf}
\caption{\textbf{Response settling after the response cue at delay 80.}
Settling was defined as the first counted response step with absolute circular
error below $20^{\circ}$; the first five transition steps were excluded. Points
denote trained models and lines denote five-model means.}
\label{fig:gaussian-settling}
\end{figure}

Together, these results show that independent Gaussian noise is useful as a
weak, non-specific degradation control. Its effects scale with perturbation
strength and maintenance duration, and are expressed most clearly as increased
decoder error, drift, and within-angle dispersion. The result establishes a
reference against which more targeted perturbations can be judged; it does not
provide evidence that Gaussian noise itself recreates psilocybin's neural or
behavioural effects.

\begin{thebibliography}{4}

\bibitem{carhartharris2014}
R.~L. Carhart-Harris \emph{et al.},
``The entropic brain: a theory of conscious states informed by neuroimaging
research with psychedelic drugs,'' Front. Hum. Neurosci. \textbf{8}, 20 (2014),
\href{https://doi.org/10.3389/fnhum.2014.00020}{doi:10.3389/fnhum.2014.00020}.

\bibitem{schartner2017}
M.~M. Schartner, R.~L. Carhart-Harris, A.~B. Barrett, A.~K. Seth, and
S.~D. Muthukumaraswamy, ``Increased spontaneous MEG signal diversity for
psychoactive doses of ketamine, LSD and psilocybin,'' Sci. Rep. \textbf{7},
46421 (2017),
\href{https://doi.org/10.1038/srep46421}{doi:10.1038/srep46421}.

\bibitem{kilpatrick2013}
Z.~P. Kilpatrick, B. Ermentrout, and B. Doiron, ``Optimizing working memory with
heterogeneity of recurrent cortical excitation,'' J. Neurosci. \textbf{33},
18999--19011 (2013),
\href{https://doi.org/10.1523/JNEUROSCI.1641-13.2013}{doi:10.1523/JNEUROSCI.1641-13.2013}.

\bibitem{ghazizadeh2021}
E. Ghazizadeh and S. Ching, ``Slow manifolds within network dynamics encode
working memory efficiently and robustly,'' PLoS Comput. Biol. \textbf{17},
e1009366 (2021),
\href{https://doi.org/10.1371/journal.pcbi.1009366}{doi:10.1371/journal.pcbi.1009366}.

\end{thebibliography}

\end{document}
